% !TeX root = ../../Book2.tex
\section{Mihlin–Hörmander 定理}
\label{b2:sec:2202}

有界符号在 \(L^2\) 上的作用已经由 Plancherel 完全描述；要得到其他指数上的界，还须控制不同频率的输出如何在空间中叠加。本节用两种条件完成这一步。Mihlin 条件逐点控制符号导数，先得到核的大小与梯度估计；Hörmander 条件只控制每个环带上的 Sobolev 范数，得到的是核差分的积分估计。两条路线都以 \(L^2\) 有界性为起点，再使用\chapref{b2:ch:18}的分解与插值方法。

\subsection{符号导数}
\label{b2:subsec:220201}

设 \(d\geq1\)，\(N\) 为非负整数，\(m\) 是定义在 \(\mathbb R^d\setminus\{0\}\) 上的复值 \(C^N\) 函数。尺度不变的导数条件写成
\[
 |\xi|^{|\alpha|}\cdot|\partial^\alpha m(\xi)|\leq A,
 \quad 0\leq|\alpha|\leq N,\quad \xi\ne0.
\]
这里 \(\alpha=(\alpha_1,\ldots,\alpha_d)\) 是多重指标，\(|\alpha|=\alpha_1+\cdots+\alpha_d\)，\(\partial^\alpha=\partial_1^{\alpha_1}\cdots\partial_d^{\alpha_d}\)。零阶条件就是 \(|m|\leq A\)。若观察半径约为 \(2^j\) 的频率环带，把它缩放到单位尺度后，一阶导数带来 \(2^j\)，而假设中的 \(|\xi|^{-1}\) 恰好抵消这个因子；更高阶导数遵循同一规律。

\ProseParagraphBreak
例如 \(m(\xi)=\xi_\ell/|\xi|\) 是零次齐次的：\(m(t\xi)=m(\xi)\)，\(t>0\)。将此式对 \(\xi\) 求导可得
\[
 (\partial^\alpha m)(t\xi)
 =t^{-|\alpha|}\cdot(\partial^\alpha m)(\xi).
\]
每个导数在单位球面上连续，因此任意固定有限阶的尺度导数界都成立。这里常数可以随阶数增长，并不要求一个常数同时控制所有阶导数。

\ProseParagraphBreak
另一个例子是 \(m(\xi)=b(\log|\xi|)\)，其中 \(b\in C^N(\mathbb R;\mathbb C)\)，并且 \(b,b',\ldots,b^{(N)}\) 全部有界。链式法则把 \(\partial^\alpha m\) 写成若干项之和，每一项都是 \(|\xi|^{-|\alpha|}\)、单位方向的有界光滑函数和某个 \(b^{(r)}(\log|\xi|)\) 的乘积，所以也满足尺度导数界。这样的符号可以沿尺度持续振荡，在原点不必有极限。乘子由几乎处处的符号决定，原点的取值不影响算子；定理要求的是离开原点后的光滑性。

\ProseParagraphBreak
尺度因子也划定了本节的范围。对 \(\alpha>0\)，\(|\xi|^{-\alpha}\) 在零频无界，因而不能是 \(L^2\) 上的有界零阶乘子；它适合放到改变可积性或正则性的分数次积分理论中。另一方面，普通的“各阶导数有界”未必包含高频处所需的衰减，不能直接代替尺度导数条件。

\subsection{二进分解}
\label{b2:subsec:220202}

先恢复\chapref{b2:ch:19}所用的分割。取实值径向函数 \(\chi\in C_c^\infty(\mathbb R^d)\)，使 \(0\leq\chi\leq1\)，在 \(|\xi|\leq1\) 时等于 \(1\)，在 \(|\xi|\geq2\) 时等于 \(0\)，并沿半径不增。令
\[
 \varphi(\xi)=\chi(\xi)-\chi(2\xi).
\]
于是 \(\varphi\geq0\)，其支撑包含于 \(\{\,1/2\leq|\xi|\leq2\,\}\)。有限和的望远镜恒等式为
\[
 \sum_{j=-M}^{N}\varphi(2^{-j}\xi)
 =\chi(2^{-N}\xi)-\chi(2^{M+1}\xi).
\]
对每个 \(\xi\ne0\)，令 \(M,N\) 趋于无穷便得到 \(\sum_{j\in\mathbb Z}\varphi(2^{-j}\xi)=1\)。每个非零频率只落在有限个相邻环带中，而且任何有限指标集 \(J\subset\mathbb Z\) 都满足
\[
 0\leq\sum_{j\in J}\varphi(2^{-j}\xi)\leq1.
\]
最后这个不等式将在取算子极限时提供共同的 \(L^2\) 控制。

\ProseParagraphBreak
把符号局部化后，令
\[
 \begin{aligned}
  m_j(\xi)&=m(\xi)\cdot\varphi(2^{-j}\xi),\\
  b_j(\zeta)&=m(2^j\zeta)\cdot\varphi(\zeta),\\
  h_j&=\mathcal F^{-1}b_j,\quad
  k_j(x)=2^{jd}\cdot h_j(2^jx).
 \end{aligned}
\]
这里第二行把所有频率环带送到同一个固定环带；Fourier 变量换元给出 \(k_j=\mathcal F^{-1}m_j\)。这样就可以先证明一个不依赖尺度的估计，再把尺度因子放回去。

\Needspace{6\baselineskip}
\begin{lemma}[固定环带上的有限导数核估计]\label{b2:lem:22-annular-kernel}
设 \(d\geq1\)，\(N\) 为非负整数，\(a\in C_c^N(\mathbb R^d;\mathbb C)\) 的支撑包含于 \(\{\,1/2\leq|\xi|\leq2\,\}\)，且
\[
 \max_{|\alpha|\leq N}\|\partial^\alpha a\|_\infty\leq M.
\]
则 \(h=\mathcal F^{-1}a\) 是光滑函数，并满足
\[
 |h(x)|+|\nabla h(x)|
 \leq C_{d,N}\cdot M(1+|x|)^{-N},
 \quad x\in\mathbb R^d.
\]
这里空间变量的一阶导数不要求符号具有 \(N+1\) 阶导数。
\end{lemma}
\begin{proof}
由于 \(a\) 支撑有界，任何多项式乘以 \(a\) 都可积。因此可以在逆 Fourier 积分内对 \(x\) 任意多次求导；这说明 \(h\) 光滑。在 \(|x|\leq1\) 时，直接估计积分给出 \(|h(x)|\leq C_dM\)。对 \(|x|>1\)，选取满足 \(|x_\ell|\geq|x|/\sqrt d\) 的坐标，分部积分 \(N\) 次：
\[
 h(x)=(-2\pi\mathrm i x_\ell)^{-N}
 \int_{\mathbb R^d}e^{2\pi\mathrm i x\cdot\xi}
 \cdot\partial_\ell^Na(\xi)\,\mathrm d\xi.
\]
边界项为零，因为 \(a\) 及所需低阶导数都紧支撑。积分区域的体积固定，于是
\[
 |h(x)|\leq C_{d,N}\cdot M|x|^{-N}.
\]
结合 \(|x|\leq1\) 的估计得到第一项的结论。

对于 \(\partial_{x_r}h\)，其频率符号是 \(2\pi\mathrm i\xi_r\cdot a(\xi)\)。在固定环带上，Leibniz 公式表明这个符号截至 \(N\) 阶的导数仍由 \(C_{d,N}M\) 控制：落在坐标 \(\xi_r\) 上的导数最多一次，落在 \(a\) 上的导数不超过 \(N\) 次。将刚才的估计用于该符号，再对 \(r=1,\ldots,d\) 求和，即得梯度界。
\end{proof}

现在对归一化符号 \(b_j\) 使用这个引理。Leibniz 公式明确给出
\[
 \partial^\alpha b_j(\zeta)
 =\sum_{\beta\leq\alpha}\binom{\alpha}{\beta}
 2^{j|\beta|}\cdot(\partial^\beta m)(2^j\zeta)
 \cdot\partial^{\alpha-\beta}\varphi(\zeta).
\]
在 \(\varphi\) 的支撑上，\(|\zeta|\) 介于 \(1/2\) 与 \(2\) 之间，因而
\[
 2^{j|\beta|}\cdot|(\partial^\beta m)(2^j\zeta)|
 \leq A|\zeta|^{-|\beta|}
 \leq 2^{|\beta|}A.
\]
所以 \(b_j\) 的 \(C^N\) 界与 \(j\) 无关。引理给出
\[
 \begin{aligned}
 |k_j(x)|&\leq C A\,2^{jd}(1+2^j|x|)^{-N},\\
 |\nabla k_j(x)|&\leq C A\,2^{j(d+1)}(1+2^j|x|)^{-N}.
 \end{aligned}
\]
本节中分割函数 \(\chi,\varphi\) 一经选定便保持不变，与它们有关的常数并入 \(C\)。

\ProseParagraphBreak
在 \(x\ne0\) 处，选整数 \(j_0\)，使 \(2^{j_0}|x|\leq1<2^{j_0+1}|x|\)。低频部分不需要衰减因子，高频部分则使用它。若 \(N>d\)，两段大小估计分别为
\[
 \begin{aligned}
 \sum_{j\leq j_0}2^{jd}
 &\leq C_d\,2^{j_0d}\leq C_d|x|^{-d},\\
 |x|^{-N}\sum_{j>j_0}2^{j(d-N)}
 &\leq C_{d,N}|x|^{-N}\cdot2^{j_0(d-N)}
 \leq C_{d,N}|x|^{-d}.
 \end{aligned}
\]
梯度估计的计算完全由同样两个级数给出，只需把其中的 \(d\) 改成 \(d+1\)，故要求 \(N>d+1\)。于是
\[
 \sum_{j\in\mathbb Z}|k_j(x)|\leq C A|x|^{-d},
 \quad
 \sum_{j\in\mathbb Z}|\nabla k_j(x)|
 \leq C A|x|^{-d-1}.
\]
在任意离开原点的紧集上，低频尾可由 \(\sum2^{jd}\) 与 \(\sum2^{j(d+1)}\) 一致控制，高频尾可由相应的衰减级数一致控制。因此核及其一阶导数的级数都局部一致收敛，
\[
 K(x)=\sum_{j\in\mathbb Z}k_j(x),\quad x\ne0,
\]
是 \(\mathbb R^d\setminus\{0\}\) 上的 \(C^1\) 函数，并满足上述大小与梯度界。这个核在离开原点处有普通函数意义，将用于给定 \(L^2\) 算子的支集外表示；其原点附近的奇性仍需由算子估计处理。

\subsection{Mihlin 乘子定理}
\label{b2:subsec:220203}

核的梯度界要求本证明选取 \(N>d+1\)。为避免在陈述中反复保留整数参数，以下采用 \(N=d+2\)。这得到\term[mihlin-chengzi-dingli]{Mihlin 乘子定理}[Mihlin multiplier theorem]的一种方便版本。

\Needspace{6\baselineskip}
\begin{theorem}[Mihlin 乘子定理的有限导数版本]\label{b2:thm:22-mihlin}
设 \(d\geq1\)，\(m\in C^{d+2}(\mathbb R^d\setminus\{0\};\mathbb C)\)，且
\[
 \sup_{\xi\ne0}|\xi|^{|\alpha|}\cdot|\partial^\alpha m(\xi)|
 \leq A<\infty,\quad |\alpha|\leq d+2.
\]
在原点任意指定 \(m\) 的值，并由
\(\widehat{T_mf}=m\cdot\widehat f\) 定义 \(L^2(\mathbb R^d;\mathbb C)\) 上的乘子。则对每个 \(1<p<\infty\)，\(T_m\) 唯一延拓为有界算子 \(L^p\to L^p\)，并且
\[
 \|T_mf\|_p\leq C_{d,p}\cdot A\|f\|_p.
\]
不同指数的延拓在相应空间的交集上一致。它还具有弱型 \((1,1)\) 延拓，弱型常数不超过 \(C_dA\)。
\end{theorem}
\begin{proof}
若 \(A=0\)，则 \(m=0\) 几乎处处，结论直接成立。以下设 \(A>0\)。零阶条件和\cref{b2:prop:22-l2}给出 \(\|T_m\|_{2\to2}\leq A\)。按前一小节构造 \(k_j\) 与 \(K\)，得到
\[
 |K(z)|\leq C A|z|^{-d},
 \quad |\nabla K(z)|\leq C A|z|^{-d-1},
 \quad z\ne0.
\]
还须证明这个 \(K\) 确实表示已经由 Fourier 变换定义的算子。令
\[
 m^{(L)}(\xi)=m(\xi)\cdot
 \sum_{|j|\leq L}\varphi(2^{-j}\xi),
 \quad K_L=\sum_{|j|\leq L}k_j.
\]
每个 \(k_j\) 都属于 \(L^1\)，故 \(T_{m^{(L)}}f=K_L*f\) 在 \(L^2\) 中成立。有限分割和介于 \(0\) 与 \(1\) 之间，且在非零频率趋于 \(1\)，于是
\[
 \|T_{m^{(L)}}f-T_mf\|_2^2
 =\int_{\mathbb R^d}|m^{(L)}-m|^2\cdot|\widehat f|^2
 \,\mathrm d\xi\longrightarrow0.
\]
这里共同控制函数是 \(4A^2\cdot|\widehat f|^2\)，可由支配收敛处理。

取紧支撑的 \(f\in L^2\)，并取紧集 \(E\subset\mathbb R^d\setminus\operatorname{supp}f\)。此时 \(f\in L^1\)，而所有差 \(x-y\)，其中 \(x\in E\)、\(y\in\operatorname{supp}f\)，组成离开原点的紧集 \(D\)。因此
\[
 \sup_{x\in E}\left|
 (K_L*f)(x)-\int K(x-y)\cdot f(y)\,\mathrm dy
 \right|
 \leq\|f\|_1\cdot\sup_{z\in D}|K_L(z)-K(z)|
 \longrightarrow0.
\]
另一方面，\(L^2\) 收敛有一个几乎处处收敛的子列。比较两个极限，得到 \(T_mf(x)=\int K(x-y)\cdot f(y)\,\mathrm dy\) 对 \(E\) 中几乎处处的 \(x\) 成立。以可数个这样的紧集穷尽支撑补集，便确认支撑外表示；该积分在支撑外绝对收敛。

把两变量核记为 \(\mathcal K(x,y)=K(x-y)\)。当 \(|x-x'|\leq|x-y|/2\) 时，连接 \(x-y\) 与 \(x'-y\) 的线段离原点至少 \(|x-y|/2\)，故微积分基本定理给出
\[
 |\mathcal K(x,y)-\mathcal K(x',y)|
 \leq C A\,\frac{|x-x'|}{|x-y|^{d+1}}.
\]
改变 \(y\) 得到另一变量的同样估计。于是 \(\mathcal K\) 满足标准 Calderón--Zygmund 核的大小与双变量光滑性条件，\(T_m\) 又已有 \(L^2\) 界。将\cref{b2:thm:18-cz-weak}和\cref{b2:thm:18-cz-lp}应用于 \(A^{-1}T_m\)，其核常数与 \(L^2\) 常数只依赖维数，得到所述弱型界与强型界。后者先在 \(L^2\cap L^p\) 上成立，再由该子空间在 \(L^p\) 中的稠密性唯一延拓；不同 \(p\) 的延拓用同一截断逼近识别。
\end{proof}

支撑外核的识别不能省略。例如常数符号 \(m=1\) 给出恒等算子，其分布核集中在原点，支撑外核却是零。一般地，支撑外核无法独自决定一个乘法项。本证明从 \(L^2\) 乘子出发，再核对它的支撑外表示，所以这种集中在对角线上的部分仍由 \(L^2\) 算子保留。

\ProseParagraphBreak
定理通常不能补上强型 \((1,1)\) 或 \(L^\infty\to L^\infty\) 两个端点。\chapref{b2:ch:18}的 Hilbert 变换给出了同时满足尺度导数条件而端点失败的例子。另一方面，若符号是支撑于一个固定环带的 \(C_c^{d+1}\) 函数，\cref{b2:lem:22-annular-kernel}已经使其核属于 \(L^1\)，Young 不等式便给出全部 \(1\leq p\leq\infty\) 的界。无限尺度之所以不同，是因为每个环带核的 \(L^1\) 范数虽一致有界，却未必对环带指标可求和。

\subsection{Hörmander 条件}
\label{b2:subsec:220204}

上面的论证把符号导数转成核的逐点衰减。若只知道导数的平方积分，仍能通过 Plancherel 得到核的带权平方积分；只要权足够强，就能控制核尾部的 \(L^1\) 积分。这条路线允许用非整数阶的光滑性，并避免要求极限核存在逐点梯度界。

\ProseParagraphBreak
本小节的 \(H^s\) 范数作用在符号的频率变量上。为免混淆，把该变量写成 \(\zeta\)，其对偶变量写成 \(x\)。对 \(s\geq0\)，采用
\[
 \|q\|_{H^s_\zeta}^2
 =\int_{\mathbb R^d}(1+4\pi^2|x|^2)^s
 \cdot|\mathcal F_\zeta q(x)|^2\,\mathrm dx.
\]
这是 \(p=2\) 的 Sobolev 范数。权关于 \(x\) 对称，因此用逆 Fourier 变换替代 \(\mathcal F_\zeta q\) 得到同一范数。

\Needspace{6\baselineskip}
\begin{lemma}[Sobolev 范数中的光滑截断]\label{b2:lem:22-sobolev-cutoff}
设 \(d\geq1\)，\(s\geq0\)，\(\psi\in C_c^\infty(\mathbb R^d;\mathbb C)\)。则乘法 \(q\mapsto\psi\cdot q\) 是 \(H^s(\mathbb R^d;\mathbb C)\) 上的有界算子，且
\[
 \|\psi\cdot q\|_{H^s}\leq C_{d,s,\psi}\cdot\|q\|_{H^s}.
\]
\end{lemma}
\begin{proof}
令 \(w_s(x)=(1+|x|)^s\)，它与 \((1+4\pi^2|x|^2)^{s/2}\) 等价。三角不等式蕴含
\[
 w_s(x)\leq w_s(x-y)\cdot w_s(y).
\]
先取 \(q\) 为 Schwartz 函数。乘积的 Fourier 变换是卷积，所以
\[
 w_s(x)\cdot|\widehat{\psi\cdot q}(x)|
 \leq
 \bigl((w_s\cdot|\widehat\psi|)*(w_s\cdot|\widehat q|)\bigr)(x).
\]
由于 \(\widehat\psi\) 是 Schwartz 函数，\(\|w_s\cdot\widehat\psi\|_1<\infty\)。Young 不等式给出
\[
 \|w_s\cdot\widehat{\psi\cdot q}\|_2
 \leq\|w_s\cdot\widehat\psi\|_1
 \cdot\|w_s\cdot\widehat q\|_2,
\]
即所需估计。

一般 \(q\in H^s\) 可在 \(H^s\) 中用 Schwartz 函数 \(q_n\) 逼近：在 Fourier 一侧先截断到有界区域，再以紧支撑光滑函数逼近，权在每个有界区域上与常数等价，逆变换即得所需 \(q_n\)。上式使 \(\psi\cdot q_n\) 在 \(H^s\) 中为 Cauchy 列；同时 \(q_n\to q\) 于 \(L^2\)，而 \(\psi\) 有界，所以 \(\psi\cdot q_n\to\psi\cdot q\) 于 \(L^2\)。两个极限一致，从而估计延续到 \(q\)。
\end{proof}

还需要把核的积分差分条件接到弱型估计上。为便于以后复用，下面单独给出卷积核版本。这个结果只要求支撑外表示和核差分积分有界，不要求核本身满足逐点大小条件。

\Needspace{6\baselineskip}
\begin{lemma}[积分 Hörmander 条件与弱型估计]\label{b2:lem:22-integral-hormander-weak}
设 \(d\geq1\)，\(T\) 是复线性有界算子 \(L^2(\mathbb R^d)\to L^2(\mathbb R^d)\)，\(\|T\|_{2\to2}\leq M\)。设 Borel 函数 \(K\in L^1_{\mathrm{loc}}(\mathbb R^d\setminus\{0\})\) 具有下列支撑外表示：对每个紧支撑的 \(f\in L^2\)，
\[
 Tf(x)=\int_{\mathbb R^d}K(x-y)\cdot f(y)\,\mathrm dy
\]
在 \(\operatorname{supp}f\) 外几乎处处成立，且右侧在这些点绝对收敛。若
\[
 H\coloneqq\sup_{v\ne0}
 \int_{|x|>2|v|}|K(x-v)-K(x)|\,\mathrm dx<\infty,
\]
则对 \(f\in L^1\cap L^2\) 和 \(t>0\)，
\[
 \bigl|\{\,x:|Tf(x)|>t\,\}\bigr|
 \leq \frac{C_d(M+H)}{t}\cdot\|f\|_1.
\]
因此 \(T\) 唯一延拓为与原 \(L^2\) 算子相容的弱型 \((1,1)\) 算子。
\end{lemma}
\begin{proof}
记 \(B=M+H\)。若 \(B=0\)，则 \(T=0\) 于 \(L^2\)，结论成立。设 \(B>0\)，并对 \(f\in L^1\cap L^2\) 在高度 \(\alpha=t/B\) 使用\cref{b2:thm:18-cz-decomposition}。该结果适用于复值函数，给出
\[
 f=g+\sum_Qb_Q,\quad
 \|g\|_2^2\leq 2^d\alpha\cdot\|f\|_1,\quad
 \sum_Q|Q|\leq\alpha^{-1}\cdot\|f\|_1,
\]
其中所选二进立方体两两不交，\(b_Q\) 支撑于 \(Q\)，\(\int b_Q=0\)，并且
\[
 \sum_Q\|b_Q\|_1\leq2\|f\|_1.
\]
每个坏块属于 \(L^2\)，坏块级数在 \(L^2\) 中收敛。这个收敛性质随后使有限坏块估计可以传到算子作用下的极限。

把每个 \(Q\) 关于其中心 \(c_Q\) 放大 \(8\sqrt d\) 倍，所得立方体记为 \(Q^*\)，并令
\[
 \Omega^*=\bigcup_QQ^*,\quad G=\mathbb R^d\setminus\Omega^*.
\]
由体积估计，
\[
 |\Omega^*|\leq C_d\cdot\frac{B}{t}\cdot\|f\|_1.
\]
好部用 \(L^2\) 有界性和 Chebyshev 不等式处理：
\[
 \bigl|\{\,|Tg|>t/2\,\}\bigr|
 \leq\frac{4M^2}{t^2}\cdot\|g\|_2^2
 \leq C_d\cdot\frac{M^2}{Bt}\cdot\|f\|_1
 \leq C_d\cdot\frac{B}{t}\cdot\|f\|_1.
\]

坏部所需的抵消正是积分条件发挥作用的地方。对 \(x\notin Q^*\)，支撑外表示和 \(\int b_Q=0\) 给出
\[
 Tb_Q(x)=\int_Q
 [K(x-y)-K(x-c_Q)]\cdot b_Q(y)\,\mathrm dy
\]
几乎处处成立。如果 \(y\in Q\) 且 \(x\notin Q^*\)，则 \(|x-c_Q|>2|y-c_Q|\)。令 \(z=x-c_Q\)、\(v=y-c_Q\)，Tonelli 定理和积分假设因此给出
\[
 \begin{aligned}
 \int_{\mathbb R^d\setminus Q^*}|Tb_Q(x)|\,\mathrm dx
 &\leq\int_Q|b_Q(y)|\cdot
   \left(\int_{\mathbb R^d\setminus Q^*}
    |K(x-y)-K(x-c_Q)|\,\mathrm dx\right)\,\mathrm dy\\
 &\leq H\|b_Q\|_1.
 \end{aligned}
\]
当 \(y=c_Q\) 时差分为零，其余点由 \(|z|>2|v|\) 的积分控制；不需为该单点另设条件。

枚举坏块并令 \(b^{(n)}\) 为前 \(n\) 项之和。在共同区域 \(G\) 上，有限和满足
\[
 \int_G|Tb^{(n)}(x)|\,\mathrm dx
 \leq H\sum_Q\|b_Q\|_1\leq2H\|f\|_1.
\]
由 \(b^{(n)}\to b=f-g\) 于 \(L^2\)，有 \(Tb^{(n)}\to Tb\) 于 \(L^2\)。选取几乎处处收敛的子列，再用 Fatou 引理，得到 \(\int_G|Tb|\leq2H\|f\|_1\)。于是
\[
 \bigl|\{\,x\in G:|Tb(x)|>t/2\,\}\bigr|
 \leq\frac{4H}{t}\cdot\|f\|_1.
\]
把这一界、好部界和 \(|\Omega^*|\) 相加，得到 \(L^1\cap L^2\) 上的弱型估计。

最后说明延拓的含义。对 \(f\in L^1\)，选有界紧支撑 \(f_n\)，使 \(\|f_n-f\|_1\leq2^{-3n}\)。对差 \(f_{n+1}-f_n\) 使用刚证的估计，可得
\[
 \sum_{n=1}^{\infty}
 \bigl|\{\,|Tf_{n+1}-Tf_n|>2^{-n}\,\}\bigr|<\infty.
\]
几乎每点最终都避开这些集合，故 \(Tf_n\) 几乎处处收敛。以其极限定义 \(Tf\)，再对 \(|Tf_n-Tf_m|\) 的分布估计令 \(m\to\infty\)，得到
\[
 \bigl|\{\,|Tf_n-Tf|>t\,\}\bigr|
 \leq\frac{C_dB}{t}\cdot\|f_n-f\|_1.
\]
同样的极限过程把弱型界传给 \(f\)，并说明不同逼近所得极限相同。用两组逼近的和与数乘可验证线性。若 \(f\in L^1\cap L^2\)，可选同时在 \(L^1,L^2\) 中逼近的截断，所得极限又与原算子的 \(L^2\) 极限相同。这给出相容性与弱型延拓的唯一性。
\end{proof}

固定 \(\eta\in C_c^\infty(\mathbb R^d\setminus\{0\})\)，使它在闭环带 \(\{\,1/2\leq|\zeta|\leq2\,\}\) 的一个邻域等于 \(1\)。对可测有界符号 \(m\)，令
\[
 B_s(m)\coloneqq\sup_{j\in\mathbb Z}
 \|m(2^j\,\cdot)\cdot\eta\|_{H^s}.
\]
固定截断 \(\eta\) 留出了一点环带边界余量，使 \(\varphi\) 和坐标乘以 \(\varphi\) 都能作为光滑乘法作用于被测量的对象。以下是\term[hormander-chengzi-dingli]{Hörmander 乘子定理}[Hörmander multiplier theorem]的环带版本。

\Needspace{6\baselineskip}
\begin{theorem}[Hörmander 环带条件]\label{b2:thm:22-hormander}
设 \(d\geq1\)，\(s>d/2\)，\(m\in L^\infty(\mathbb R^d;\mathbb C)\)。固定 \(\eta\in C_c^\infty(\mathbb R^d\setminus\{0\})\)，使它在 \(\{\,1/2\leq|\xi|\leq2\,\}\) 的邻域等于 \(1\)，并假设
\[
 B_s(m)=\sup_{j\in\mathbb Z}
 \|m(2^j\,\cdot)\cdot\eta\|_{H^s(\mathbb R^d)}<\infty.
\]
则 \(L^2\) 乘子 \(T_m\) 对每个 \(1<p<\infty\) 唯一延拓为有界算子 \(L^p\to L^p\)，且
\[
 \|T_mf\|_p\leq C_{d,p,s,\eta}
 \bigl(\|m\|_\infty+B_s(m)\bigr)\cdot\|f\|_p.
\]
它还具有与 \(L^2\) 定义相容的弱型 \((1,1)\) 延拓，弱型常数不超过
\(C_{d,s,\eta}\bigl(\|m\|_\infty+B_s(m)\bigr)\)。
\end{theorem}
\begin{proof}
仍用本节固定的 \(\varphi\)，并记
\[
 q_j(\zeta)=m(2^j\zeta)\cdot\eta(\zeta),\quad
 b_j=\varphi\cdot q_j,\quad h_j=\mathcal F^{-1}b_j,\quad
 k_j(x)=2^{jd}\cdot h_j(2^jx).
\]
因为 \(\eta=1\) 于 \(\operatorname{supp}\varphi\) 的邻域，\(b_j=m(2^j\,\cdot)\cdot\varphi\)，与二进分解中的符号相同。\cref{b2:lem:22-sobolev-cutoff}分别用于截断 \(\varphi\) 和 \(2\pi\mathrm i\zeta_r\cdot\varphi(\zeta)\)，给出
\[
 \|b_j\|_{H^s}+\sum_{r=1}^{d}
 \|2\pi\mathrm i\zeta_r\cdot b_j\|_{H^s}
 \leq C_{d,s,\eta}\cdot B_s(m).
\]
Plancherel 与 Sobolev 范数的定义因此说明
\[
 \|(1+|x|)^s\cdot h_j\|_2
 +\|(1+|x|)^s\cdot\nabla h_j\|_2
 \leq C\cdot B_s(m).
\]
这些 \(h_j\) 是光滑函数：\(b_j\) 有界且紧支撑，任何多项式乘以 \(b_j\) 都属于 \(L^1\)，可以在逆变换中对空间变量求导。光滑性来自有限频带上的可积性与空间变量求导。

令 \(\gamma=s-d/2>0\)。由 Cauchy--Schwarz 不等式，
\[
 \|h_j\|_1+\|\nabla h_j\|_1
 \leq C\cdot B_s(m)\cdot
 \left(\int_{\mathbb R^d}(1+|x|)^{-2s}\,\mathrm dx\right)^{1/2}
 \leq C_{d,s,\eta}\cdot B_s(m).
\]
更精确地，对 \(R\geq1\)，极坐标给出
\[
 \int_{|x|>R}(1+|x|)^{-2s}\,\mathrm dx
 \leq C_d\int_R^\infty r^{d-1-2s}\,\mathrm dr
 =\frac{C_d}{2s-d}\cdot R^{d-2s},
\]
所以
\[
 \int_{|x|>R}|h_j(x)|\,\mathrm dx
 \leq C_{d,s,\eta}\cdot B_s(m)R^{-\gamma}.
\]
这一步把带权 \(L^2\) 信息变成了两个不同用途的 \(L^1\) 信息：梯度的全空间积分控制小位移，核的尾积分控制大位移。

取有限指标集 \(J\subset\mathbb Z\)，令
\[
 K_J=\sum_{j\in J}k_j,\quad
 m_J(\xi)=m(\xi)\cdot\sum_{j\in J}\varphi(2^{-j}\xi),
 \quad T_Jf=K_J*f.
\]
每个 \(k_j\in L^1\)，故 \(K_J\in L^1\)，卷积与 \(L^2\) 乘子 \(T_{m_J}\) 一致。分割的非负性保证 \(|m_J|\leq\|m\|_\infty\)，因而
\[
 \|T_J\|_{2\to2}\leq\|m\|_\infty
\]
与有限集 \(J\) 无关。下一步证明积分 Hörmander 常数也与 \(J\) 无关。

固定 \(v\ne0\)。若 \(2^j|v|\leq1\)，由
\[
 k_j(x-v)-k_j(x)
 =-\int_0^1v\cdot\nabla k_j(x-tv)\,\mathrm dt
\]
及平移不变性可得
\[
 \begin{aligned}
 \int_{|x|>2|v|}|k_j(x-v)-k_j(x)|\,\mathrm dx
 &\leq|v|\cdot\|\nabla k_j\|_1\\
 &=2^j|v|\cdot\|\nabla h_j\|_1
 \leq C\cdot B_s(m)\cdot2^j|v|.
 \end{aligned}
\]
若 \(2^j|v|>1\)，则在 \(|x|>2|v|\) 上同时有 \(|x|>|v|\) 与 \(|x-v|>|v|\)。分别在两个积分中换元，使用上面的尾积分界，得到
\[
 \begin{aligned}
 &\int_{|x|>2|v|}|k_j(x-v)-k_j(x)|\,\mathrm dx\\
 &\quad\leq
 \int_{|z|>2^j|v|}|h_j(z)|\,\mathrm dz+
 \int_{|z|>2^j|v|}|h_j(z)|\,\mathrm dz
 \leq C\cdot B_s(m)\cdot(2^j|v|)^{-\gamma}.
 \end{aligned}
\]
选择 \(j_0\) 满足 \(2^{j_0}|v|\leq1<2^{j_0+1}|v|\)。两部分级数满足
\[
 \sum_{j\leq j_0}2^j|v|\leq2,\quad
 \sum_{j>j_0}(2^j|v|)^{-\gamma}
 \leq\frac{1}{1-2^{-\gamma}}.
\]
对任意有限 \(J\) 取子和不会增大这些上界，所以
\[
 \sup_{v\ne0}\int_{|x|>2|v|}
 |K_J(x-v)-K_J(x)|\,\mathrm dx
 \leq C_{d,s,\eta}B_s(m).
\]

有限和核是光滑的 \(L^1\) 函数，满足所需支撑外表示。于是\cref{b2:lem:22-integral-hormander-weak}给出所有 \(T_J\) 的一致弱型界。其 \(L^2\) 伴随的核为
\[
 K_J^*(x)=\overline{K_J(-x)}.
\]
变量替换 \(x\mapsto-x\)、\(v\mapsto-v\) 表明 \(K_J^*\) 具有同一个积分条件，伴随的 \(L^2\) 范数也相同，所以 \(T_J^*\) 具有同样的弱型界。记 \(D=\|m\|_\infty+B_s(m)\)；若 \(D=0\) 则结论平凡，否则先对 \(D^{-1}T_J\) 使用\cref{b2:thm:17-marcinkiewicz}，由弱型 \((1,1)\) 和强型 \((2,2)\) 得到 \(1<p<2\) 的界。对 \(D^{-1}T_J^*\) 作相同插值，再由 \(p'\in(1,2)\) 的对偶性得到 \(2<p<\infty\) 的界。因此
\[
 \|T_Jf\|_p\leq C_{d,p,s,\eta}\cdot D\|f\|_p,
 \quad1<p<\infty,
\]
常数与 \(J\) 的大小和位置无关。

最后令 \(J=\{\,-L,\ldots,L\,\}\)。分割恒等式和 \(|m_J|\leq\|m\|_\infty\) 使 \(T_Jf\to T_mf\) 于 \(L^2\)。对 \(f\in L^p\cap L^2\)，选取几乎处处收敛的子列，再用 Fatou 引理，将统一 \(L^p\) 界传到 \(T_mf\)。在 \(L^1\cap L^2\) 上也可把弱型分布估计传到极限：先把阈值 \(t\) 降为任意 \(t'<t\)，使用几乎处处收敛和 Fatou，再令 \(t'\uparrow t\)。强型延拓由 \(L^p\cap L^2\) 的稠密性得到，弱型延拓由前一引理证明末尾的逼近步骤得到；所有这些延拓与原 \(L^2\) 算子相容。
\end{proof}

严格不等式 \(s>d/2\) 在两处承担同一任务：使逆权 \((1+|x|)^{-s}\) 属于 \(L^2\)，并使尾积分指数 \(\gamma=s-d/2\) 为正。在 \(s=d/2\) 时，当前尾积分计算出现对数发散，高频环带级数也不再衰减。这指出了本证明不能到达端点的原因；证明失效本身并不构成某个附加条件下端点结论的反例。

\ProseParagraphBreak
有限环带和的估计通过 \(L^2\) 极限传到原算子。与 Mihlin 证明相比，保留下来的是足以估计坏块的核差分积分，而不是逐点梯度。这正是 Sobolev 环带条件可以降低光滑性要求的原因。

\ProseParagraphBreak
还可用平移乘子看清“充分条件”的含义。若 \(a\in\mathbb R^d\)，则 \(m(\xi)=e^{2\pi\mathrm i a\cdot\xi}\) 给出 \(T_mf(x)=f(x+a)\)，它在每个 \(L^p\) 上等距。但当 \(a\ne0\) 时，至少一个一阶导数的模恒为非零常数，乘上 \(|\xi|\) 后无界，所以不满足 Mihlin 条件。Mihlin 条件的作用在于直接检验符号，并使所得控制在尺度分解下保持一致。

\ProseParagraphBreak
最后给出下一节所需的一族乘子版本。标量算子范数一致有界并不能单独推出先在位置上求平方和的估计；此处统一的 Mihlin 条件还提供统一核结构，因而可以调用\cref{b2:prop:18-hilbert-valued-cz}。

\Needspace{6\baselineskip}
\begin{proposition}[统一 Mihlin 条件下的平方和估计]\label{b2:prop:22-uniform-mihlin-square}
设 \(d\geq1\)，\(1<p<\infty\)，\(I\) 是有限或可数指标集。对每个 \(\nu\in I\)，设 \(m_\nu\in C^{d+2}(\mathbb R^d\setminus\{0\};\mathbb C)\)，并假设
\[
 \sup_{\nu\in I}\sup_{\xi\ne0}
 |\xi|^{|\alpha|}\cdot|\partial^\alpha m_\nu(\xi)|
 \leq A<\infty,\quad|\alpha|\leq d+2.
\]
以各自 \(L^2\) 定义的相容 \(L^p\) 延拓记 \(T_{m_\nu}\)。则对任意复值函数族 \((f_\nu)_{\nu\in I}\)，只要右侧有限，就有
\[
 \left\|\left(\sum_{\nu\in I}|T_{m_\nu}f_\nu|^2\right)^{1/2}\right\|_p
 \leq C_{d,p}\cdot A
 \left\|\left(\sum_{\nu\in I}|f_\nu|^2\right)^{1/2}\right\|_p.
\]
常数与指标集的大小无关。
\end{proposition}
\begin{proof}
先设 \(I=\{1,\ldots,n\}\)，并令 \(H=\mathbb C^n\)，赋予 Euclidean 范数。在 \(L^2(\mathbb R^d;H)\) 上定义对角算子
\[
 \mathbb T(f_1,\ldots,f_n)
 =(T_{m_1}f_1,\ldots,T_{m_n}f_n).
\]
逐坐标 Plancherel 给出
\[
 \|\mathbb Tf\|_{L^2(H)}^2
 =\sum_{\nu=1}^{n}\int|m_\nu(\xi)|^2
 \cdot|\widehat f_\nu(\xi)|^2\,\mathrm d\xi
 \leq A^2\cdot\|f\|_{L^2(H)}^2.
\]
对每个符号，用本节二进分解构造支撑外核 \(K_\nu\)。全部符号共享常数 \(A\)，因此核的大小、梯度和支撑外表示也共享同一组常数。令
\[
 \mathbb K(z)=\operatorname{diag}
 \bigl(K_1(z),\ldots,K_n(z)\bigr),\quad z\ne0.
\]
对角矩阵在 \(H\) 上的算子范数等于对角元模的最大值，故
\[
 \|\mathbb K(z)\|_{H\to H}
 =\max_{1\leq\nu\leq n}|K_\nu(z)|
 \leq C A|z|^{-d}.
\]
当 \(|h|\leq|z|/2\) 时，各坐标的梯度界还给出
\[
 \|\mathbb K(z+h)-\mathbb K(z)\|_{H\to H}
 \leq C A\,\frac{|h|}{|z|^{d+1}}.
\]
对两变量核 \(\mathbb K(x-y)\)，改变任一变量都满足相应差分估计。支撑外表示由有限个坐标的表示逐项组合；此处有限维 Bochner 积分就是坐标积分。

这样，\(\mathbb T\) 满足\cref{b2:prop:18-hilbert-valued-cz}的全部条件：有 \(L^2(H)\) 起点、算子范数意义的标准核估计，以及紧支撑输入的支撑外表示。该命题的常数不依赖 Hilbert 空间的维数。应用它并对 \(A\) 归一化，即得
\[
 \|\mathbb Tf\|_{L^p(H)}
 \leq C_{d,p}\cdot A\|f\|_{L^p(H)},
\]
这正是有限族的平方和估计。先对 \(L^2(H)\cap L^p(H)\) 输入应用，再由稠密性推广，得到陈述中的相容 \(L^p\) 算子。

若 \(I\) 可数，枚举其元素并对前 \(n\) 项使用已证估计。右侧的有限平方和由完整平方和控制；左侧随着 \(n\) 增大逐点单调增加。对其 \(p\) 次方使用单调收敛定理，得到可数族结论。
\end{proof}

这个向量估计同时控制一组频带上的乘子。将它与\chapref{b2:ch:19}的平方函数定理结合，就能把逐频带的估计重新合成为原函数的范数；下一节将据此识别 Bessel 势范数。
